% Laget automatisk av elevark/lag_elevark.py fra sim/magnetfelt-ledere.html. Endringer her blir overskrevet.
\documentclass[11pt]{article}
\usepackage{elevark}
\begin{document}
\elevarkhode{Fysikk 2 · Magnetiske felt og krefter}{Magnetfelt rundt ledere}{https://lektorbo.no/sim/magnetfelt-ledere.html}{Simuleringen}

Rundt en lang, rett leder med strøm er det et magnetfelt med sirkulære feltlinjer. Her står lederne vinkelrett på skjermen, så strømmen går ut av $(⊙)$ eller inn i $(⊗)$ planet. Dra lederne og målepunktet $P$ rundt, og se hvordan feltene legger seg sammen.

\section*{Oppgaver}

\begin{oppgaveliste}
\oppgave \fase{gjett}{Gjett}Velg «Én leder». Strømmen går ut av planet. Går feltlinjene med eller mot klokka? Bruk høyrehåndsregelen, og sjekk med kompassnålene. Hva skjer når du snur strømmen?
\svar
\oppgave \fase{test}{Regn}Strømmen er 10~A. Regn ut magnetfeltet 5,0~cm fra lederen, og sjekk med $P$. I Norge er jordas magnetfelt omtrent 50~µT. Ville et kompass merket lederen?
\svar
\oppgave \fase{gjett}{Gjett}Hvor sterkt er feltet 2,5~cm fra lederen? Sammenlign med det elektriske feltet rundt en punktladning. Avtar de to feltene på samme måte med avstanden?
\svar
\oppgave \fase{test}{Regn}Velg «To ledere, samme retning». Hva er feltet midt mellom lederne? Regn ut kraften per meter mellom lederne, og avgjør om de tiltrekker eller frastøter hverandre. Slå på «Krefter på lederne» og sjekk.
\svar
\oppgave \fase{test}{Regn}Velg «To ledere, motsatt retning». Regn ut feltet midt mellom lederne. Hvorfor blir det dobbelt så sterkt som fra én leder?
\svar
\oppgave \fase{gjett}{Gjett}Velg «10~A og 5~A». Strømmene går i samme retning, og det er 12~cm mellom lederne. Hvor mellom lederne er feltet null? Regn først, og finn punktet med $P$.
\svar
\oppgave \fase{utfordring}{Utfordring}Forklar med høyrehåndsreglene hvorfor to ledere med strøm i samme retning tiltrekker hverandre. Bruk feltet fra den ene lederen og kraften $F = IlB$ på den andre.
\svar
\end{oppgaveliste}

\end{document}
